На множестве $\Omega$ стандартным образом определяется размерность Хаусдорфа (см. например \cite[Секция 6]{Edgar}).
Для непустого подмножества $F\subset \R^n$ и $s > 0$ определим $s$-мерную меру Хаусдорфа множества $F$ следующим образом:
$$\mathcal H^s(F) := \lim_{\delta\to0} \inf \left\{\sum_{i=1}^\infty \left({\rm diam} \ U_i \right)^s \ : \ F\subset \bigcup_{i=1}^\infty  U_i, \  0\leqslant {\rm diam} \ U_i \leqslant \delta \right\}.$$

Размерность Хаусдорфа множества  $F\subset \R^n$ определяется по формуле:
$${\rm dim}_H F := \inf\{ s > 0 \ : \ \mathcal H^s(F)=0\}.$$



Мы приведём определение самоподобных подмножеств множества $\Omega$ (см., например, \cite{falconer2003fractal}).

\begin{definition}
Множество $E\subset\Omega$ называется самоподобным, если существуют $m\in\N$,
$m\geqslant2$, $0< r_1, \dots, r_m<1$ и функции $f_j : \Omega \to \Omega$, $j=1,\dots, m$ такие, что
$$\rho(f_j(x), f_j(y)) = r_j \rho(x,y), \ \forall \ x,y \in \Omega, \ j=1,\dots, m$$
и $E=\bigcup_{j=1}^m f_j(E).$
\end{definition}

